% 汇盈探店 · 探店员小程序使用手册（xelatex 编译）
\documentclass[12pt,a4paper,fontset=none]{ctexart}
\newcommand{\manualtitle}{探店员小程序使用手册}
% 共享排版样式：管理员手册 / 探店员手册 通用前言（xelatex + ctex 编译）
\usepackage[margin=2.2cm,headheight=14pt]{geometry}
\usepackage{xcolor}
\usepackage{hyperref}
\usepackage{booktabs}
\usepackage{tabularx}
\usepackage{array}
\usepackage{enumitem}
\usepackage{fancyhdr}
\usepackage{needspace}
\usepackage{colortbl}

\hypersetup{colorlinks=true,linkcolor=black,urlcolor=blue,pdftitle=\manualtitle}

% 字体（macOS 中文字体）
\setCJKmainfont{Songti SC}[BoldFont={Songti SC Bold}]
\setCJKsansfont{Hiragino Sans GB}[BoldFont={Hiragino Sans GB W6}]
\setCJKmonofont{Menlo}
\setmainfont{Times New Roman}

% 品牌色
\definecolor{brand}{RGB}{242,95,66}
\definecolor{brandlight}{RGB}{255,240,236}
\definecolor{graytext}{RGB}{138,141,146}

% 标题样式：黑体 + 品牌色
\ctexset{
  section/format={\Large\bfseries\sffamily\color{brand}},
  subsection/format={\large\bfseries\sffamily\color{black}},
  subsubsection/format={\normalsize\bfseries\sffamily},
}

% 页眉页脚
\pagestyle{fancy}
\fancyhf{}
\fancyhead[L]{\small\sffamily\color{graytext}\manualtitle}
\fancyhead[R]{\small\sffamily\color{graytext}\thepage}
\renewcommand{\headrulewidth}{0.4pt}
\renewcommand{\headrule}{\color{brand!60}\hrule width\headwidth height 0.4pt}

% 提示框：\tipbox{标题}{内容}
\newcommand{\tipbox}[2]{%
  \par\vspace{0.5em}\noindent
  \fcolorbox{brand}{brandlight}{\begin{minipage}{\dimexpr\linewidth-2\fboxsep-2\fboxrule\relax}
    \textbf{\color{brand}#1}\par\smallskip #2
  \end{minipage}}\par\vspace{0.5em}}

% 警示框（红色）：\warnbox{内容}
\newcommand{\warnbox}[1]{%
  \par\vspace{0.5em}\noindent
  \fcolorbox{red!70}{red!8}{\begin{minipage}{\dimexpr\linewidth-2\fboxsep-2\fboxrule\relax}
    \textbf{\color{red!70!black}注意}\par\smallskip #1
  \end{minipage}}\par\vspace{0.5em}}

% 步骤列表：紧凑
\setlist[enumerate]{itemsep=2pt,topsep=4pt,leftmargin=2.2em}
\setlist[itemize]{itemsep=2pt,topsep=4pt,leftmargin=1.8em}

% 表格默认样式
\renewcommand{\arraystretch}{1.25}

% 封面
\newcommand{\manualcover}[4]{%
  \begin{titlepage}
    \centering
    \vspace*{2.5cm}
    {\color{brand}\rule{6cm}{2.5pt}}\\[1.2cm]
    {\Huge\bfseries\sffamily #1\par}
    \vspace{0.8cm}
    {\LARGE\sffamily #2\par}
    \vspace{2.5cm}
    {\large\color{graytext} 汇盈探店 · 门店服务品质暗访系统\par}
    \vspace{0.4cm}
    {\large\color{graytext} #3\par}
    \vfill
    {\color{brand}\rule{6cm}{2.5pt}}\\[0.6cm]
    {\normalsize\color{graytext} #4\par}
  \end{titlepage}}


\begin{document}

\manualcover{探店员小程序使用手册}{汇盈探店 · 微信小程序}{版本 1.0 · 2026 年 8 月}{本手册面向探店员，覆盖从登录、查看任务、到店打卡、填写问卷、上传小票到提交的全部流程。}

\tableofcontents
\newpage

\section{手册说明}

\subsection{适用对象}
本手册面向「汇盈探店」微信小程序的探店员。你的账号由管理方在后台创建（或经审核批准），登录后即可查看指派给自己的探店任务。

\subsection{核心概念}
\begin{itemize}
\item \textbf{探店任务}：管理员分派给你的一次到店检查。每名探店员的任务相互独立，问卷只属于你本人。
\item \textbf{打卡}：开始任务前，系统会重新获取你的位置，判断是否到达门店；通过后才能开始填写问卷。
\item \textbf{得分 / 不得分}：问卷每题二选一。选择「不得分」时必须填写问题说明，并按题目要求上传照片或录音证据。
\item \textbf{消费小票}：部分任务要求上传小票。上传后可用识别功能辅助录入菜品，识别结果需逐项核对。
\item \textbf{闪光点}：提交前必填的一项总结性描述。
\end{itemize}

\subsection{快速上手}
\begin{enumerate}
\item 登录小程序，进入「首页」。
\item 在「下一项任务」卡点击进入任务详情，查看要求。
\item 到达门店后点击「开始探店」，完成定位打卡。
\item 依次完成环境卫生、为客服务、产品品质三类题目（50 道）。
\item 在「产品品质」分类点击「菜品记录」，上传小票并填写菜品。
\item 填写闪光点，确认提交。提交后问卷锁定，不能回看或修改。
\end{enumerate}

\section{登录}

\subsection{登录流程}
打开小程序，在登录页输入「用户名」与「密码」（可点密码框右侧眼睛图标切换显示/隐藏），点击「登录」。

\begin{itemize}
\item 账号或密码错误：提示「用户名或密码不正确」。
\item 管理员账号同样可以登录本小程序（界面与探店员一致，只是通常没有分配任务）。
\item 连续多次失败会触发限流，提示「操作过于频繁，请稍后重试」，稍等一分钟再试。
\item 登录成功后进入定位说明页；若账号还在审核中、未通过或已停用，会看到对应状态说明页。
\end{itemize}

\warnbox{当前版本已隐藏「注册 / 忘记密码」入口。新账号由管理方在后台创建，密码忘记时请联系管理方重置。}

\subsection{密码规则（自助设置时）}
密码需 \textbf{8–15 个字符，且必须包含大写字母和英文符号}（如 A、\$、! 等）。管理员在后台创建的账号不受此限制，但同样建议设置 8 位以上。

\section{首页}

首页自上而下包含：

\begin{itemize}
\item \textbf{下一项任务卡}：系统按「进行中优先，其次待开始，时间早的在前」选出下一项任务；点击整卡进入任务详情。进行中的任务按钮显示「继续填写」，待开始显示「查看任务」。
\item \textbf{执行提醒}：保持普通顾客身份；到店后重新定位。
\item \textbf{统计}：待开始、进行中、逾期三个数字。
\end{itemize}

无任务时显示「当前没有待执行任务——新任务发布后会显示在这里」，可点「查看全部日程」。加载失败时点「重新加载」。

\section{日程}

日程页「我的日程」按日期分组列出你的全部任务，组标题为任务开始日期（如「8月11日 周二」）。每张任务卡显示状态、类型（暗访/明察）、门店、任务名与时间段，点击进入任务详情。

\begin{itemize}
\item 顶部显示任务总数。
\item 下拉页面可强制刷新最新任务。
\item 跨天任务额外显示「截止 ×月×日」。
\end{itemize}

\section{问卷入口页}

问卷页列出与问卷相关的任务：

\begin{itemize}
\item \textbf{可执行任务}：待开始的任务按钮为「查看并开始」，进行中的为「继续填写」，点击进入任务详情。
\item \textbf{等待处理}：已逾期的任务仅展示，备注「等待管理员延期」——需要管理员延期后才能继续。
\item \textbf{已完成}：已提交的任务仅展示提交时间。
\item \textbf{明察任务}：显示「暂未启用」说明，第一阶段暂不开放明察问卷。
\end{itemize}

\section{我的}

「我的」页包含：

\begin{itemize}
\item \textbf{用户信息卡}：头像（姓名首字）、姓名、角色「探店员」与账号状态（已启用 / 申请正在审核中 / 申请未通过 / 已停用）。
\item \textbf{定位与隐私}：查看位置用途与隐私说明（位置仅用于判断是否到店打卡，不持续记录轨迹）。
\item \textbf{填写帮助}：查看答题、证据、菜品填写说明。
\item \textbf{退出登录}：确认后退出，下次需重新输入账号密码。
\end{itemize}

\section{任务详情与开始探店}

任务详情页展示：

\begin{itemize}
\item 头部卡：状态、类型、门店、地址、时间窗。
\item 「本次要求」：问卷版本、预算上限（如「不超过 500 元」）、任务状态。
\item 「菜品要求」：必点菜品与建议菜品清单。
\item 特别说明（如有）。
\item 已提交任务显示「提交后不能回看答案、菜品、说明或证据」；逾期任务显示警示与「返回日程」。
\end{itemize}

\subsection{开始探店（五步）}
\begin{enumerate}
\item 点击「开始探店」，系统开始获取你的当前位置（提示「正在获取现场位置…」）。
\item 系统核对是否到达门店（「正在核对是否到达门店…」）。
\item 打卡通过后创建问卷（「位置已通过，正在创建问卷…」）。
\item 打开问卷（「问卷已准备好，正在打开…」）。
\item 若进行中已有草稿，按钮显示「继续填写」，直接回到问卷。
\end{enumerate}

\tipbox{定位权限}{系统在开始时才会重新定位，不会持续记录轨迹。若拒绝定位授权，仍可查看任务说明，但不能开始探店；可在弹窗中「打开设置」重新授权。}

\subsection{打卡失败的几种结果}

\begin{center}
\begin{tabular}{p{5.2cm}p{8cm}}
\toprule
提示 & 含义与处理 \\
\midrule
现场位置已通过 & 打卡成功，可开始填写 \\
当前距离门店较远，请到店后重试 & 位置不在门店打卡范围内，走近门店后重试 \\
定位结果已过期，请重新定位 & 定位超过 60 秒，重新定位 \\
定位精度不足，请移动到较开阔位置后重试 & 室内信号差，到开阔处重试 \\
当前位置未通过，请重新定位 & 定位异常，重新定位 \\
暂时无法确认位置，请稍后重试 & 稍后重试；多次失败可联系管理方 \\
\bottomrule
\end{tabular}
\end{center}

\subsection{开始时的常见拦截}
\begin{itemize}
\item 「还没到探店开始时间，请稍后再试」——任务未到开始时间，稍后再来。
\item 「任务已逾期，请等待管理员延期后刷新任务」——任务逾期，联系管理方延期。
\item 「请先完成现场打卡后再开始」——打卡凭证已过期，需重新打卡。
\end{itemize}

\section{填写问卷}

问卷页顶部为三个分类页签：\textbf{环境卫生（1–20 题）、为客服务（21–42 题）、产品品质（43–50 题）}，状态行显示「已完成 N / 50」。

\subsection{作答}
每题选择「得分」或「不得分」：
\begin{itemize}
\item \textbf{得分}：保存成功即计该题满分。
\item \textbf{不得分}：进入不得分页，必须填写「问题说明」，并按题目配置上传证据。
\item 不得分改判为得分：需进入改判确认确认——确认后清除原说明与证据；取消则保持原答案不变。
\end{itemize}

\tipbox{保存机制}{答案会实时保存。状态行显示「正在保存… / 已保存」；若提示「保存失败，请重试」，说明网络异常，点击重试；若提示「内容已更新，请刷新后重试」，按提示重新加载页面。}

\subsection{问题说明与证据}
在不得分页：
\begin{itemize}
\item \textbf{问题说明}：必填，描述现场看到的问题（支持输入法语音输入）。
\item \textbf{问题照片}（题目要求时）：点击「拍照或相册（最多 N 张）」，N 为该题配置的照片上限（默认 3 张）；已上传照片显示「已保存照片」。
\item \textbf{问题音频}（题目要求时）：点击「录制音频（最长 N 秒）」，N 为该题配置的时长上限（默认 10 秒），到时长自动停止；每题最多 3 段，单段至少 1 秒。
\item 底部「确认」提交该题；「取消」放弃本次修改，原说明与证据保持不变。
\item 单个文件不能超过 10 MB。
\end{itemize}

\section{菜品记录与消费小票}

完成全部 50 题后，在「产品品质」分类底部点击「菜品记录」。

\subsection{上传消费小票}
\begin{itemize}
\item 点击「＋ 添加小票」拍照或从相册选择；最多 \textbf{9 张}，仅支持 图片，每张不超过 \textbf{10 MB}。
\item 任务要求小票时，至少完成 1 张小票上传才能提交。
\item 上传失败的条目可「重试」；已上传失败可「重新选择」。
\item 删除小票时需选择处理方式：「同时移除未修改的识别菜品」或「保留菜品，转为手工记录」。
\end{itemize}

\subsection{小票识别辅助录入}
\begin{itemize}
\item 点击「识别并带入菜品」，系统识别小票并带入候选菜品。
\item 识别状态：「正在排队识别」「正在识别小票」「已识别 N 道菜」「已带入 N 道菜，请核对标记内容」「暂时无法识别这张小票，您仍可以手动填写菜品」。
\item \textbf{识别结果只是候选}，请逐项核对；带「请核对」「系统推算」「已人工修改」「来自小票」标记的内容需要特别留意。
\item 识别失败不影响填写，可完全手工录入。
\end{itemize}

\subsection{填写菜品}
每道菜需确认：
\begin{itemize}
\item \textbf{实际菜名}（售罄/未点时留空）。
\item \textbf{状态}：已点 / 已替代（需填替代原因）/ 已售罄（需填售罄原因）/ 未点（仅建议菜品）；必点菜品不能选择「未点」。
\item \textbf{数量}：1–999 整数；单价与行总价选填（金额最多两位小数）。
\item \textbf{评价}：记录上菜时间、外观、温度、味道、份量、口感等现场观察。
\item \textbf{满意度}：整单「满意 / 未满意」二选一。
\item 任务指定的必点菜品不可删除；可点「＋ 点击添加其他菜品」补充额外菜品。
\end{itemize}

\section{其他记录与提交}

\subsection{其他记录}
\begin{itemize}
\item \textbf{闪光点}：必填，描述本次体验中值得肯定之处。
\item \textbf{其他建议}：选填，填写问卷未列出的改进建议。
\end{itemize}

\subsection{提交}
\begin{enumerate}
\item 点击「确定」，弹出确认框：「提交后将无法回看或修改本次问卷，请确认所有内容已经完成。」
\item 点击「确认提交」。提交中请勿关闭页面。
\item 提交成功页显示提交时间；本次问卷锁定，不能回看或修改。如需说明，请联系管理方。
\end{enumerate}

\warnbox{提交前请确认：50 题全部作答、不得分题说明与证据齐全、闪光点已填写、任务要求的小票已上传。提交不可撤销。}

\subsection{提交失败的几种情况}

\begin{center}
\begin{tabular}{p{5.4cm}p{8cm}}
\toprule
提示 & 含义与处理 \\
\midrule
问卷已过最终截止时间，请等待管理员延期后刷新任务 & 超过问卷截止时间，联系管理方延期 \\
请完成全部问卷题目后再提交 & 有题目未作答 \\
不得分题目缺少必要证据 & 不得分题的照片/录音不完整 \\
请填写本次探店亮点 & 闪光点未填写 \\
请至少上传一张已完成确认的消费小票 & 任务要求的小票未上传成功 \\
任务已逾期，请等待管理员延期后刷新任务 & 任务逾期，联系管理方延期 \\
\bottomrule
\end{tabular}
\end{center}

\section{常见问题}

\begin{center}
\begin{tabular}{p{5.4cm}p{8cm}}
\toprule
现象 & 处理 \\
\midrule
提示「暂时无法连接服务」 & 检查网络后重试；仍失败可联系管理方 \\
提示「登录已失效」 & 重新登录 \\
打卡一直失败 & 确认定位权限已开启、位置在门店范围内；仍失败联系管理方申请打卡豁免 \\
任务显示「已逾期」 & 联系管理方延期；延期后任务恢复 \\
提示「还没到探店开始时间」 & 任务尚未到开始时间，到点再来 \\
提示「内容已更新，请刷新后重试」 & 下拉刷新页面后重试 \\
照片/录音上传失败 & 检查文件大小（≤10MB）与格式，重试 \\
\bottomrule
\end{tabular}
\end{center}

\appendix

\section{附录 A：任务状态对照}

\begin{center}
\begin{tabular}{ll}
\toprule
状态 & 中文显示 \\
\midrule
assigned & 待开始 \\
in\_progress & 进行中 \\
overdue & 已逾期 \\
submitted & 已提交 \\
cancelled & 已取消 \\
\bottomrule
\end{tabular}
\end{center}

\section{附录 B：填写帮助原文}

小程序「我的 → 填写帮助」中的五项说明：

\begin{enumerate}
\item \textbf{开始任务}：在任务详情页点击「开始探店」，系统会重新获取当前位置并确认是否在门店范围内；通过后才能开始填写问卷。
\item \textbf{答题规则}：50 道题按环境卫生、为客服务、产品品质三部分完成。每题选择「得分」或「不得分」；选择不得分时必须填写问题说明，并按提示上传照片或录音。
\item \textbf{菜品与消费小票}：按要求填写实际点到的菜品和评价。任务要求消费小票时，需上传至少 1 张成功的原图；上传后可识别小票辅助填写菜品，识别结果请务必逐项核对。
\item \textbf{提交前检查}：闪光点为必填项。提交后无法回看或修改，请确认全部内容后再提交。
\item \textbf{定位与隐私}：位置仅用于判断是否到店打卡，不会持续记录行动轨迹。拒绝授权时可以查看任务说明，但不能开始探店。
\end{enumerate}

\end{document}
